\section{Introduction}
\label{sec:introduction-v132}

A degeneracy locus usually forgets the map that defines it.  We ask
whether its nonreduced scheme structure can recover that map, without
an ambient embedding or a marking of the normal directions.  For the
multiplication-failure schemes of quadratic pencils the answer is
exact: there is one uniform infinitesimal order at which the entire
pencil first becomes visible.  The inverse theorem applies to every
pencil, including singular pencils.  Its regular locus then carries
the spectral and critical correspondences studied below.

\subsection{The principal theorem}
Let \(V\) be a complex vector space of dimension \(n\ge3\), and set
\[
 W=\Sym^2V,\qquad N=\binom{n+1}{2},\qquad
 R\in\Gr(2,W),\qquad S_R=W/R,\qquad p=N-2.
\]
The algebra \(B_R=\C\oplus V\oplus S_R\) has multiplication
\(\Sym^2V\twoheadrightarrow S_R\), with \((V\oplus S_R)S_R=0\).
On \(X_R=\Gr(n,V\oplus S_R)\) let \(\mathcal K\) denote the
tautological bundle.  For any \(m\ge2\), define
\begin{equation}\label{eq:failure-intro}
 \widehat D_R=
 V\!\left(
 \Fitt_0\coker[
 \Sym^m(\OO\oplus\mathcal K)\longrightarrow
 B_R\otimes\OO]\right).
\end{equation}
This is the literal maximal-minor scheme of the multiplication map;
it is independent of \(m\ge2\).  Its reduction is the fixed
Schubert divisor where \(\mathcal K\to V\) drops rank.  Let
\(B=\Gr(n,S_R)\), let \(\mathfrak b_R\) be its ideal in
\(\widehat D_R\), and put
\[
 Z_R^{[k]}=V(\mathfrak b_R^{k+1})\subset\widehat D_R,
 \qquad d=n+2p=n^2+2n-4.
\]
The neighbourhood is taken along the whole projective socle
Grassmannian, not at one framed affine point.

\begin{theorem}[Sharp intrinsic finite reconstruction]
\label{thm:principal-finite-v143}
For fixed \(n\ge3\), the following assertions hold on all of
\(\Gr(2,\Sym^2V)\).

\smallskip
\noindent\textup{(i)} For every pair of pencils \(R,R'\),
\[
 Z_R^{[d]}\simeq Z_{R'}^{[d]}
 \quad\Longleftrightarrow\quad
 R'=gR\quad\hbox{for some }g\in\PGL(V).
\]
The isomorphism of neighbourhoods is abstract and unmarked.
For each \(0\le k<d\), the isomorphism class of \(Z_R^{[k]}\)
is independent of \(R\).  Thus \(d\) is the smallest uniform
reconstruction order.

\smallskip
\noindent\textup{(ii)} In fixed tensor coordinates, the first relation
spaces define a closed immersion
\[
 \Gr(2,W)\longrightarrow
 \Gr\!\left(\binom N2,\Sym^d\Hom(U,V)^*\right),\qquad \dim U=n.
\]
This assertion includes infinitesimal parameter schemes and singular
pencils.

\smallskip
\noindent\textup{(iii)} These relations glue to the actual universal
order-\(d\) neighbourhood over the relative socle Grassmannian.
It is finite flat there of rank
\[
                  \binom{n^2+d}{d}-\binom N2,
\]
projective and flat over the pencil parameter scheme, and compatible
with arbitrary complex base change with the relative socle retained.
\end{theorem}

Theorem~\ref{thm:sharp-finite-pencil} proves \textup{(i)};
Theorems~\ref{thm:finite-parameter-immersion} and
\ref{thm:universal-finite-neighbourhood} prove \textup{(ii)--(iii)}.
The distinction between their inputs is essential: \textup{(i)} uses
an abstract global scheme, whereas \textup{(ii)--(iii)} retain the
relative socle and a tensor presentation.  No equivalence of unmarked
moduli stacks is asserted.

Here is the mechanism behind the theorem.  In graph coordinates
\(T:\mathcal U\to V\), where \(\mathcal U\) is tautological on
\(B\), the homogeneous ideal is
\[
          (\det T)I_p(\gamma_R\Sym^2T),
          \qquad \gamma_R:W\twoheadrightarrow S_R.
\]
It is generated in degree \(d\).  The abstract finite scheme recovers
\(B\), its degree-one conormal, and the first relation module.
The reduced determinant cone has two rank-one rulings.  Only its
\(\Pj(V)\)-ruling is projectively trivial over \(B\); a
Schubert line distinguishes the other one.  Regular descent therefore
recovers a single constant projective coordinate transformation on
\(V\), rather than unrelated changes for different coefficients.
Cancelling the determinant line then recovers the residual coefficient
module.  For pencils its exterior representation is irreducible;
exterior duality identifies its coefficient line with the full
Pluecker line of \(R\).  The first visible degree alone is formal;
these intrinsic identifications are the substantive inverse argument.

\subsection{Spectral and critical correspondences}
On the regular locus, the reconstructed pencil determines its torsion
spectral sheaf, not merely its determinant.  Theorems
\ref{thm:spectral-finite-torelli} and
\ref{thm:relative-spectral-readout} recover both the elementary
divisors and their scheme-theoretic Fitting incidence loci.  Theorem
\ref{thm:fixed-spectral-classification} applies this to every fixed
discriminant and reduced-rank stratum: the classical dominance
specializations are realized by flat families of the finite
neighbourhoods, and every strict change first becomes distinguishable
at order \(d\).  The qualifier \emph{reduced} matters; the full
spectral Fitting schemes need not remain fixed.

There is a further functorial output.  The reciprocal critical
correspondence is the zero scheme of the two score equations on the
invertible members of a pencil, with the data matrix allowed to vary.
Theorem~\ref{thm:critical-correspondence-v143} recovers its isomorphism
class from the finite failure invariant and constructs it in families.
Thus the passage from finite failure to likelihood has a specified
algebraic intermediate object, rather than a parallel numerical
analogy.  For a supplied definite real realization,
Theorem~\ref{thm:real-critical-cover-v143} gives a congruence-covariant
open family of data, a finite etale critical algebra of rank \(2r-3\),
and globally labelled real analytic critical sections.  Its residue
and separated-root argument builds on the complete proof of
Theorem~\ref{thm:totally-real-reciprocal-likelihood}.  That latter
theorem proves the assertion formulated in
\cite[Conjecture~4.5]{FevolaMandelshtamSturmfels}, for unrestricted
real data and arbitrary eigenvalue multiplicities.  Neither a real
form nor a definite member is inferred from an unmarked complex
scheme, and positivity of the data or of all critical points is not
asserted.

\subsection{Proof dependencies and further results}
The principal argument is independent of exterior recombination.
The following table specifies the logical inputs; the implication
from reconstructed matrices to the critical correspondence is formal,
whereas its totally real open family requires an additional proof.

\begin{center}
\small
\begin{tabular}{@{}p{.22\linewidth}p{.33\linewidth}p{.37\linewidth}@{}}
\toprule
Step & Input & Output \\
\midrule
Intrinsic coefficients & Literal Fitting ideal; oriented rank-one
rulings & One projective frame and the full Pluecker line \\[3pt]
Finite reconstruction & Coefficient inverse; single relation degree &
Sharp uniform order; parameter immersion \\[3pt]
Relative spectral readout & Glued finite neighbourhood; regular pencil &
Spectral sheaf and Fitting incidence schemes \\[3pt]
Critical correspondence & Reconstructed pencil; universal score equations &
Congruence-covariant algebraic correspondence \\[3pt]
Totally real cover & Supplied real definite realization; residues and signs &
Finite etale algebra and labelled real critical sections \\
\bottomrule
\end{tabular}
\end{center}

The supplement retains every web, contraction, polarized, polar,
and boundary theorem with its full proof.  In particular,
Theorem~\ref{thm:main-web-reconstruction} recovers all basepoint-free
four-dimensional webs with smooth Jacobian quartic; the all-ranks
contraction theorem and exact Jacobian--Casimir factorization supply
its recombination mechanism.  They are substantial further results,
not hidden hypotheses of the pencil theorem.  The first-layer K3
construction, hyperelliptic interpretation, component-fibre
classification, and primary boundary calculations remain available
at their original level of strength.

\subsection{Relation to classical results}
The Weierstrass--Segre classification, its orbit closures, the harmonic
decomposition, and exterior-support bounds are classical inputs,
identified at the statements where they are used.  In particular,
\cite[Theorem~1.1 and Corollary~2.1]{FevolaMandelshtamSturmfels}
classify regular pencils by their labelled elementary divisors;
\cite[equation (2)]{TrevisiolSymmetric} gives the selfadjoint orthogonal orbit
calculation underlying the fixed-spectrum discussion.  The inverse
problem here starts instead with an unmarked failure scheme.
The exact differential dictionary for the auxiliary contraction is
given in Section~\ref{sec:operator-dictionary-v143}; it separates
classical harmonic tools from the normalization of the specified map.

Ballico's paper \cite{Ballico93} is a directly relevant predecessor.
Its complete text has not been obtained for the present comparison.
Consequently we do not infer nonanticipation from its title, first
page, or later citations.  The six theorem-level comparisons requiring
that text are recorded explicitly in the accompanying referee response.
This documentary limitation is separate from the hypotheses and proofs
of the theorems stated here.
